
% DEFAULT PACKAGE SETUP
\usepackage{setspace,graphicx,epstopdf,amsmath,amsfonts,amssymb,amsthm,versionPO,threeparttable}
\usepackage{marginnote,datetime,enumitem,subfig,rotating,fancyvrb,tabularx,booktabs}
\usepackage{float}
% \usepackage[longnamesfirst]{natbib}
\usepackage[]{natbib}
\newcolumntype{Y}{>{\raggedleft\arraybackslash}X}% raggedleft column X
\interfootnotelinepenalty=10000 %% Completely prevent breaking of footnotes

% Define \numberthis
\newcommand{\numberthis}{\addtocounter{equation}{1}\tag{\theequation}}

\usdate
\usepackage{lscape}
% \usepackage{endfloat}
% \usepackage[nolists]{endfloat}
\usepackage{rotating} 

\usepackage[justification=justified]{caption}
% \usepackage{subcaption}
% These next lines allow including or excluding different versions of text
% using versionPO.sty

\excludeversion{notes}  % Include notes?
\includeversion{links}  % Turn hyperlinks on?

\usepackage{eurosym}
% Turn off hyperlinking if links is excluded
\iflinks{}{\hypersetup{draft=true}}

% Notes options
\ifnotes{%
 \usepackage[margin=1in,paperwidth=10in,right=2.5in]{geometry}%
 \usepackage[textwidth=1.4in,shadow,colorinlistoftodos]{todonotes}%
}{%
 \usepackage[margin=1in]{geometry}%
 \usepackage[disable]{todonotes}%
}

\DeclareMathOperator*{\argmin}{arg\,min}


% Allow todonotes inside footnotes without blowing up LaTeX
% Next command works but now notes can overlap. Instead, we'll define
% a special footnote note command that performs this redefinition.
% \renewcommand{\marginpar}{\marginnote}%

% Save original definition of \marginpar
\let\oldmarginpar\marginpar

% Workaround for todonotes problem with natbib (To Do list title comes out wrong)
\makeatletter\let\chapter\@undefined\makeatother % Undefine \chapter for todonotes

% Define note commands
\newcommand{\smalltodo}[2][] {\todo[caption={#2}, size=\scriptsize, fancyline, #1] {\begin{spacing}{.5}#2\end{spacing}}}
\newcommand{\rhs}[2][]{\smalltodo[color=green!30,#1]{{\bf RS:} #2}}
\newcommand{\rhsnolist}[2][]{\smalltodo[nolist,color=green!30,#1]{{\bf RS:} #2}}
\newcommand{\rhsfn}[2][]{% To be used in footnotes (and in floats)
 \renewcommand{\marginpar}{\marginnote}%
 \smalltodo[color=green!30,#1]{{\bf RS:} #2}%
 \renewcommand{\marginpar}{\oldmarginpar}}
% \newcommand{\textnote}[1]{\ifnotes{{\noindent\color{red}#1}}{}}
\newcommand{\textnote}[1]{\ifnotes{{\colorbox{yellow}{{\color{red}#1}}}}{}}

% Command to start a new page, starting on odd-numbered page if twoside option
% is selected above
\newcommand{\clearRHS}{\clearpage\thispagestyle{empty}\cleardoublepage\thispagestyle{plain}}

% Number paragraphs and subparagraphs and include them in TOC
\setcounter{tocdepth}{2}

% JF-specific includes:

\usepackage{indentfirst} % Indent first sentence of a new section.
\usepackage{endnotes} % Use endnotes instead of footnotes
\usepackage{rfs}  % JF-specific formatting of sections, etc.
% \usepackage[labelfont=bf,labelsep=period]{caption} % Format figure captions
\captionsetup[table]{labelsep=none}

\usepackage{color}%Color
% Define theorem-like commands and a few random function names.
\newtheorem{condition}{CONDITION}
\newtheorem{corollary}{COROLLARY}
\newtheorem{proposition}{PROPOSITION}
\newtheorem{obs}{OBSERVATION}
\newtheorem{lem}{LEMMA}
\newtheorem{theorem}{THEOREM}
\newtheorem{assumption}{ASSUMPTION}
\newcommand{\argmax}{\mathop{\rm arg\,max}}
\newcommand{\sign}{\mathop{\rm sign}}
\newcommand{\defeq}{\stackrel{\rm def}{=}}

\usepackage{multirow}
\usepackage{booktabs}


\usepackage{amsmath}
\usepackage{amsthm}
\usepackage{amsfonts}
\usepackage{relsize}
\usepackage{tikz}
\usepackage{pgfplotstable} 
\usepgfplotslibrary{dateplot}
\usetikzlibrary{fit}
\usepackage{graphicx}
\usepackage{eurosym}
\usepackage{standalone}
% \usepackage{filecontents}

\usepackage[section]{placeins} % keep figures in the same section


\usepackage{appendix}
\appendixtitleon
\newcommand\independent{\protect\mathpalette{\protect\independenT}{\perp}}
\def\independenT#1#2{\mathrel{\rlap{$#1#2$}\mkern2mu{#1#2}}}

%\onehalfspacing
%\linespread{1.25}

\setstretch{1.5}
%\doublespacing

\definecolor{zahi}{RGB}{255,0,0}
\newcommand{\zahi}[1]{\textcolor{zahi}{[zahi: #1]}}



%%%%% Pgfplots %%%%%
\usepackage{pgfplots}
\usepgfplotslibrary{groupplots}

\pgfplotsset{compat=newest}
\usetikzlibrary{patterns}

%%%% Table caption package %%%%%
\usepackage{ragged2e}
\usepackage{longtable}

%%%% Figure imports %%%%
\usepackage{import}

%%%% Color cells in tables %%%%
\usepackage{tabulary,colortbl}
\newcolumntype{C}[1]{>{\centering\let\newline\\\arraybackslash\hspace{0pt}}m{#1}}


\definecolor{andrea}{RGB}{255,0,0}
\newcommand{\andrea}[1]{\textcolor{andrea}{[andrea: #1]}}

\definecolor{yang}{RGB}{255,0,0}
\newcommand{\yang}[1]{\textcolor{yang}{[yang: #1]}}
\newcommand{\jay}[1]{\textcolor{yang}{[j: #1]}}


\usepackage[colorlinks=true,linkcolor=red,anchorcolor=black,citecolor=black, filecolor=black,menucolor=black,runcolor=black,urlcolor=black]{hyperref}
\usepackage{parskip}



%%%%%%%%%%%%%%%%% For plots: get min, max etc


\newcommand*{\GetElement}[4]{%
    \pgfplotstablegetelem{#2}{#3}\of{#1}%
    \let#4\pgfplotsretval
}

% get the first element of a column and store it in a macro
% #1: table
% #2: column (name or [index]<index>)
% #3: macro for value
\newcommand*{\GetFirstElement}[3]{%
    \GetElement{#1}{0}{#2}{#3}%
}

% get the last element of a column and store it in a macro
% #1: table
% #2: column (name or [index]<index>)
% #3: macro for value
\newcommand*{\GetLastElement}[3]{%
    \pgfplotstablegetrowsof{#1}%
    \pgfmathsetmacro\lastrowidx{int(\pgfplotsretval - 1)}%
    \GetElement{#1}{\lastrowidx}{#2}{#3}%
}

% get maximum and minimum value of a column and store them in macros
% #1: table
% #2: column
% #3: macro for max
% #4: macro for min
\newcommand*{\GetMinMax}[4]{%
    \pgfmathsetmacro#3{-16383}%
    \pgfmathsetmacro#4{16383}%
    \pgfplotstableforeachcolumnelement{#2}\of{#1}\as\cellValue{%
        \ifx\cellValue\@empty\else
            \pgfmathsetmacro{#3}{max(#3,\cellValue)}%
            \pgfmathsetmacro{#4}{min(#4,\cellValue)}%
        \fi
    }
}

% get average of column and store it in a macro
% #1: table
% #2: column
% #3: macro for average
\newcommand*{\GetAverage}[3]{%
    \pgfplotstablegetrowsof{#1}%
    \let\rownumber\pgfplotsretval
    \def#3{0}%
    \pgfplotstableforeachcolumnelement{#2}\of{#1}\as\cellValue{%
        \ifx\cellValue\@empty
            % don't count emplty cells
            \pgfmathsetmacro{\rownumber}{\rownumber - 1}%
        \else
            \pgfmathsetmacro{#3}{#3 + \cellValue}%
        \fi
    }
    \pgfmathsetmacro{#3}{#3 / \rownumber}%
}

%%%%%%%%%%%%%%%%%%%%%%%%%%%%%%%%%%%%%%%%%%%%%%%%%%%%%%

%%%% Quintile colors (Green to Red) %%%% 
\definecolor{q1}{rgb}{ .973,  .412,  .42}
\definecolor{q2}{rgb}{ .992,  .749,  .486}
\definecolor{q3}{rgb}{ 1,  .922,  .518}
\definecolor{q4}{rgb}{ .678,  .827,  .498}
\definecolor{q5}{rgb}{ .388,  .745,  .482}
%%%%%%%%%%%%%%%%%%%%%%%%%%%%%%%%%%%%%%%%